\chapter{Belajar Menyelesaikan: Neural Operator dan Simulator}
\label{ch:operator}

Bab ini juga tidak berasal dari satu mata kuliah. Kuliah geometric deep learning memberi bahasa graf dan
ekuivariansi, kuliah machine learning lanjut memberi transformasi Fourier dan neural ODE, dan dosennya
sendiri, Johannes Brandstetter, termasuk peneliti yang membentuk bidang ini, dari message passing PDE
solver sampai transformer untuk fisika. Sisanya datang dari bacaan.

PINN mempelajari satu solusi. Kalau syarat awalnya berubah sedikit, pelatihan harus diulang. Neural
operator dan simulator yang dipelajari mengambil tujuan yang lebih ambisius: belajar \emph{cara}
menyelesaikan sebuah keluarga persamaan, sehingga solusi untuk syarat awal, parameter, atau geometri baru
bisa diperoleh dalam satu lintasan maju, sering ribuan kali lebih cepat dari solver klasik.

\section{Dari fungsi ke fungsi}

Solver numerik pada dasarnya adalah sebuah \istilah{operator}: ia memetakan satu fungsi, misalnya medan
temperatur awal atau distribusi konduktivitas, ke fungsi lain, yaitu solusinya. Jaringan saraf biasa
memetakan vektor berukuran tetap ke vektor berukuran tetap, sehingga terikat pada satu resolusi grid.
\istilah{Operator learning} ingin belajar pemetaan antar ruang fungsi itu sendiri, sehingga model yang
dilatih pada grid kasar bisa dievaluasi pada grid yang lebih halus \citep{kovachki2023}.

\subsection{DeepONet}

DeepONet \citep{lu2021} bertumpu pada teorema aproksimasi universal untuk operator. Operator $G$ yang
memetakan fungsi input $v$ ke solusi $G(v)$, dievaluasi di titik $y$, didekati sebagai
\begin{equation}
G(v)(y)\approx\sum_{k=1}^pb_k(v)\,t_k(y).
\end{equation}
Jaringan \istilah{branch} menerima nilai $v$ di sejumlah titik sensor tetap dan mengeluarkan koefisien
$b_k$. Jaringan \istilah{trunk} menerima koordinat $y$ dan mengeluarkan fungsi basis $t_k(y)$. Bentuk ini
mengingatkan pada ekspansi deret seperti Fourier atau POD (\cref{ch:penemuan}), dengan perbedaan bahwa
basis dan koefisiennya sama-sama dipelajari. Karena trunk menerima koordinat kontinu, keluarannya bisa
dievaluasi di titik mana pun.

\subsection{Fourier neural operator}

Fourier neural operator \citep{li2021} berangkat dari bentuk solusi banyak PDE linear: konvolusi syarat
awal dengan sebuah kernel, seperti kernel Gauss pada persamaan panas di \cref{ch:mlat}. Satu lapisan FNO
memakai operator integral dengan kernel yang dipelajari,
\begin{equation}
(\mathcal Kv)(x)=\int\kappa(x-y)\,v(y)\dd y.
\end{equation}
Karena kernel hanya bergantung pada selisih $x-y$, operator ini adalah konvolusi, dan teorema konvolusi
(\cref{ch:cv}) mengubahnya menjadi perkalian di domain frekuensi:
\begin{equation}
\mathcal Kv=\mathcal F^{-1}\big(\bm R\cdot\mathcal Fv\big),
\end{equation}
dengan $\bm R$ bobot yang dipelajari untuk setiap frekuensi. FNO hanya menyimpan bobot untuk sejumlah
kecil frekuensi terendah dan membuang sisanya. Satu lapisan lengkap menjumlahkan operator spektral ini
dengan transformasi linear lokal, lalu menerapkan nonlinearitas. Dengan FFT, biayanya sebanding dengan
$N\log N$ untuk $N$ titik grid, dan karena bobotnya didefinisikan per frekuensi, bukan per titik grid,
model yang sama bisa dijalankan pada resolusi yang berbeda. Pemotongan frekuensi juga punya sisi buruk:
struktur yang tajam, seperti gelombang kejut, sulit diwakili oleh sedikit mode Fourier.

\section{Simulator berbasis graf}

Grid teratur jarang ditemui dalam simulasi teknik. Mesh menyesuaikan diri dengan geometri, dan partikel
bergerak bebas. Di sini jaringan graf (\cref{ch:gdl}) menjadi alat alami. Graph Network-based Simulator
\citep{sanchez2020} memperlakukan setiap partikel sebagai simpul dan menghubungkan partikel yang
berdekatan dengan sisi. Arsitekturnya mengikuti pola \istilah{encode--process--decode}: keadaan setiap
partikel dikodekan menjadi vektor laten, beberapa langkah message passing menyebarkan informasi
interaksi, lalu decoder mengeluarkan percepatan setiap partikel. Posisi baru dihitung dengan integrator
sederhana, dan langkah itu diulang. Model yang dilatih pada air, pasir, dan cairan kental ini
menghasilkan rollout yang stabil untuk ribuan langkah.

MeshGraphNets \citep{pfaff2021} membawa gagasan yang sama ke mesh simulasi, dengan dua jenis sisi:
sisi mesh yang mengikuti konektivitas elemen, dan sisi ``dunia'' antara simpul yang berdekatan di ruang
walau tidak tersambung di mesh, misalnya untuk kain yang menyentuh dirinya sendiri. Model ini bahkan
belajar menyesuaikan resolusi mesh. Message Passing Neural PDE Solvers \citep{brandstetter2022} merumuskan
message passing sebagai solver PDE umum dan menunjukkan bahwa ia mencakup skema numerik klasik seperti
beda hingga dan volume hingga sebagai kasus khusus. Hubungan ini sudah kita lihat di \cref{ch:numerik}:
volume hingga adalah message passing fluks.

\section{Masalah rollout}

Hampir semua simulator yang dipelajari bekerja secara autoregresif: tebak langkah berikutnya, lalu
pakai tebakan itu sebagai input untuk langkah selanjutnya. Galat kecil di setiap langkah tidak hanya
menumpuk, tetapi juga bisa diperbesar. Misalkan galat di langkah $t$ adalah $e_t$, model menambahkan galat
baru $\delta$ di setiap langkah, dan kepekaan langkah model terhadap inputnya diwakili faktor $J$. Maka
$e_{t+1}\approx Je_t+\delta$, dan setelah $T$ langkah
\begin{equation}
e_T\approx\delta\sum_{k=0}^{T-1}J^k.
\end{equation}
Kalau $|J|<1$, galat jenuh di sekitar $\delta/(1-|J|)$. Kalau $|J|$ sedikit lebih besar dari satu,
galat tumbuh eksponensial. Dengan $J=1{,}01$, setelah 500 langkah faktor $J^{500}\approx145$. Model yang
tampak sempurna dalam uji satu langkah bisa meledak dalam rollout panjang.

Ada masalah yang lebih halus. Selama pelatihan, model hanya melihat input dari simulasi sejati. Selama
rollout, ia melihat input dari tebakannya sendiri, yang sedikit berbeda distribusinya. Ini covariate
shift yang sama dengan behavioral cloning di \cref{ch:rl}. Beberapa obat sudah dicoba: menambahkan noise
pada input selama pelatihan agar model terbiasa dengan input yang tidak sempurna \citep{sanchez2020},
melatih dengan rollout beberapa langkah, atau, seperti PDE-Refiner \citep{lippe2023}, memperhalus
setiap tebakan dengan beberapa langkah denoising mirip model difusi, supaya komponen frekuensi tinggi
dengan amplitudo kecil tidak diabaikan. Analisis spektrum sering menunjukkan bahwa model kehilangan
detail halus lebih dulu, dan detail itulah yang kemudian merusak stabilitas.

\section{Model fondasi untuk fisika dan cuaca}

Keberhasilan paling mencolok terjadi di prakiraan cuaca, di mana data reanalisis global berpuluh tahun
tersedia. FourCastNet \citep{pathak2022} memakai operator Fourier yang adaptif. GraphCast \citep{lam2023}
memakai message passing di atas mesh ikosahedral berjenjang yang membungkus bumi, dan dalam banyak ukuran
mengungguli sistem prakiraan numerik operasional untuk sepuluh hari ke depan, dengan waktu hitung kurang
dari satu menit. GenCast \citep{price2025} memakai model difusi untuk menghasilkan ensemble prakiraan
yang probabilistik. Aurora \citep{bodnar2024} dilatih pada beragam data atmosfer sebagai model fondasi yang
bisa disetel ke tugas-tugas baru, seperti kualitas udara.

Gagasan model fondasi juga dibawa ke PDE secara umum. Pretraining pada banyak sistem fisika sekaligus
\citep{mccabe2023} ternyata membantu ketika model disetel ke sistem baru dengan data sedikit. Kumpulan
data besar seperti The Well \citep{ohana2024} menyediakan simulasi dari banyak bidang fisika untuk
melatih dan menguji model seperti ini. Universal Physics Transformers \citep{alkin2024} memisahkan model
dari diskretisasi: titik mesh atau partikel dikodekan ke ruang laten berukuran tetap, dinamika berjalan di
ruang laten, dan decoder bisa ditanya di titik mana pun. \citet{gupta2022} membandingkan arsitektur
secara sistematis pada berbagai skala ruang dan waktu.

\section{Pertanyaan yang tetap harus diajukan}

Kembali ke tiga pertanyaan di \cref{ch:peta}. Fisika masuk ke model-model ini terutama lewat data dan
arsitektur: simulasi yang benar sebagai data latih, dan ekuivariansi, lokalitas, serta konservasi sebagai
struktur jaringan. Di luar jangkauan data latih, misalnya pada bilangan Reynolds lain atau geometri yang
jauh berbeda, jaminannya tipis, dan pengujian harus dilakukan secara eksplisit. Biaya pelatihannya
besar, sering setara ribuan simulasi klasik untuk menghasilkan data, sehingga keuntungan baru muncul
kalau model dipakai berkali-kali, misalnya untuk optimasi desain, kendali waktu nyata, atau ensemble
besar. Ulasan \citet{lino2023} dan \citet{vinuesa2022} memberi gambaran yang seimbang tentang kekuatan
dan batas pendekatan ini dalam mekanika fluida.

\sumberbab{Bab ini dirangkai dari kuliah Deep Learning: Geometric Techniques, Machine Learning: Advanced
Techniques, dan bacaan selama tugas akhir. Makalah utamanya: \citet{lu2021}, \citet{li2021},
\citet{kovachki2023}, \citet{sanchez2020}, \citet{pfaff2021}, \citet{brandstetter2022}, \citet{lippe2023},
dan \citet{lam2023}.}
